\subsection{ASYMPTOTIC EXPANSIONS WITH VARIABLE COEFFICIENTS}

One can generalize the concept of principal part, and that of an asymptotic expansion, in the following way. Let E be a comparison scale formed of real (resp. complex) functions such that, for each of them, there is a set in F on which the function does not vanish at any point. Further, let C be a set of functions in \(\mathcal{H}(\mathfrak{F}, V)\) satisfying the following three conditions:

\((\mathrm{CO}_{\mathrm{I}})\) For every function \(\mathbf{a} \in \mathcal{C}\) one has \(\mathbf{a} \preccurlyeq 1\) .

\((\mathrm{CO}_{\mathrm{II}})\) The relation \(\mathbf{a} \ll 1\) for a function \(\mathbf{a} \in \mathcal{C}\) implies \(\mathbf{a} = 0\) .

\((\mathrm{CO}_{\mathrm{III}})\mathcal{C}\) is a vector space over R (resp. C).

Now let \(\mathbf{f}\) be any function in \(\mathcal{H}(\mathfrak{F}, V)\) ; if there exist a function \(g \in \mathcal{E}\) and a nonzero function \(\mathbf{a} \in \mathcal{C}\) such that \(\mathbf{f} - \mathbf{a}g \ll g\) one will say that \(\mathbf{a}g\) is a principal part of \(\mathbf{f}\) , relative to the comparison scale \(\mathcal{E}\) and to the domain of coefficients \(\mathcal{C}\) . If such a principal part exists, it is unique: for suppose that there are two principal parts \(\mathbf{a}_1g_1\) and \(\mathbf{a}_2g_2\) ; one cannot have \(g_1 \ll g_2\) since from \((\mathrm{CO_I})\) one could deduce that \(\mathbf{a}_1g_1 \ll g_2\) and \(\mathbf{f} - \mathbf{a}_1g_1 \ll g_1 \ll g_2\) , so \(\mathbf{f} \ll g_2\) ; but one would also have \(\mathbf{a}_2g_2 \ll g_2\) and consequently \(\mathbf{a}_2 \ll 1\) contradicting the hypothesis \(\mathbf{a}_2 \neq 0\) and \((\mathrm{CO_{II}})\) . So one must have \(g_1 = g_2\) ; from the relations \(\mathbf{f} - \mathbf{a}_1g_1 \ll g_1, \mathbf{f} - \mathbf{a}_2g_1 \ll g_1\) one deduces that \((\mathbf{a}_2 - \mathbf{a}_1)g_1 \ll g_1\) whence \(\mathbf{a}_2 - \mathbf{a}_1 \ll 1\) , and consequently \(\mathbf{a}_2 = \mathbf{a}_1\) , by \((\mathrm{CO_{II}})\) and \((\mathrm{CO_{III}})\) .

Example. For \(x\) real tending to \(+\infty\) the periodic bounded functions on \(\mathbf{R}\) , having the same period \(\tau\) , satisfy conditions \((\mathrm{CO}_{\mathrm{I}})\) , \((\mathrm{CO}_{\mathrm{II}})\) and \((\mathrm{CO}_{\mathrm{III}})\) : if \(\lim_{x\to +\infty}a(x) = 0\) then for every \(\varepsilon >0\) there is an \(x_0\) such that \(|a(x)|\leqslant \varepsilon\) for every \(x\geqslant x_0\) ; one deduces that \(|a(x)|\leqslant \varepsilon\) for \(0\leqslant x\leqslant \tau\) too, since there exists an integer \(n\) such that \(x + n\tau \geqslant x_0\) , and such that \(a(x) = a(x + n\tau)\) ; since \(\varepsilon\) is arbitrary one has \(a(x) = 0\) on \([0,\tau]\) , hence everywhere.

With the notation of V, p.~222, we will say that \(\sum_{\lambda \leqslant \alpha} \mathbf{a}_{\lambda} g_{\lambda}\) , where the \(\mathbf{a}_{\lambda}\) belong to \(\mathcal{C}\) and all but a finite number of them are zero, is an asymptotic expansion of \(\mathbf{f}\) with coefficients in \(\mathcal{C}\) , to precision \(g_{\alpha}\) , if \(\mathbf{f} - \sum_{\lambda \leqslant \alpha} \mathbf{a}_{\lambda} g_{\lambda} \ll g_{\alpha}\) ; for every index \(\mu\) such that \(\mathbf{a}_{\mu} \neq 0\) then \(\mathbf{a}_{\mu} g_{\mu}\) is the principal part of \(\mathbf{f} - \sum_{\lambda < \mu} \mathbf{a}_{\lambda} g_{\lambda}\) , relative to \(\mathcal{E}\) and to \(\mathcal{C}\) , which proves the uniqueness of the asymptotic expansion of \(\mathbf{f}\) (to precision \(g_{\alpha}\) ) when it exists.

The methods given in \(n^{\circ}\) 3 (V, p.~223) for forming an asymptotic expansion for \(f_{1} + f_{2}\) or for \([f_{1}.f_{2}]\) , starting from given asymptotic expansions for \(f_{1}\) and \(f_{2}\) , can again be applied to expansions with variable coefficients, so long as the \([a_{\lambda}.b_{\mu}]\) belong to the domain of coefficients C corresponding to the normed space V or admit an asymptotic expansion with coefficients in C.

